% Generated from locale/tr/content/sets-functions-relations/size-of-sets/introduction.tex; do not hand-edit.


\olfileid[tr]{sfr}{siz}{int}

\olsection{Giriş}

Georg Cantor 1870'lerde küme teorisini geliştirirken amaçlarından biri, Orta
Çağ düşünürlerinin deyişiyle edimsel bir sonsuzluk olan sonsuz bir topluluk
fikrini kabul edilebilir kılmaktı. Bunun önemli bir bölümünü, farklı kümelerin
\emph{büyüklüğünü} ele alış biçimi oluşturuyordu. $a$, $b$ ve $c$ birbirinden
farklıysa $\{a, b, c\}$ kümesi sezgisel olarak $\{a, b\}$ kümesinden
\emph{daha büyüktür}. Peki sonsuz kümeler ne olacak? Bunların hepsi aynı
büyüklükte midir? Öyle olmadıkları ortaya çıkar.

Buradaki ilk önemli fikir numaralandırmadır. Her sonlu kümeyi, bütün
!!{element}s{}ını listeleyerek sıralayabiliriz. Bazı sonsuz kümelerin de,
listenin kendisinin sonsuz olmasına izin verirsek bütün !!{element}s{}ını
listeleyebiliriz. Böyle kümelere !!{enumerable} denir. Cantor'un bu bölümün
sonunda tümüyle anlayacağımız şaşırtıcı sonucu, bazı sonsuz kümelerin
!!{enumerable} olmadığıdır.

